\documentclass[11pt]{article}

% ------------------------------------------------------------------
% PACKAGES
% ------------------------------------------------------------------
\usepackage[margin=1in]{geometry}
\usepackage{graphicx}
\usepackage{booktabs}
\usepackage{amsmath, amssymb}
\usepackage{hyperref}
\usepackage{algorithm}
\usepackage{algpseudocode}
\usepackage{siunitx}
\usepackage{caption}
\usepackage{subcaption}
\usepackage{xcolor}

\hypersetup{
    colorlinks=true,
    linkcolor=blue,
    citecolor=blue,
    urlcolor=blue
}

\title{\textbf{Fair Scheduling for Multi-Tenant SSDs: \\ Design, Implementation, and Evaluation Using MQSim}}
\author{Qiyue Chen}
\date{December 2025}

\begin{document}

\maketitle

% ------------------------------------------------------------------
% ABSTRACT
% ------------------------------------------------------------------
\begin{abstract}
Solid-state drives (SSDs) increasingly serve multi-tenant environments where 
several applications issue concurrent I/O requests. When tenants share internal
SSD hardware resources, default round-robin scheduling often leads to unfair
performance allocation, starvation of smaller tenants, and violation of quality-of-service (QoS) guarantees.
This project uses the MQSim SSD simulator to implement and evaluate several 
fairness-aware scheduling algorithms—Deficit Round Robin (DRR), Quick Fair 
Queueing (QFQ), and a novel Min--Max fairness scheduler—inside MQSim's 
Transaction Scheduling Unit (TSU). We develop multi-flow workloads to study 
throughput fairness, latency isolation, and weighted service proportionality.
Our experiments (placeholder results in this report) demonstrate that RR 
exhibits large fairness violations, DRR improves equal sharing, QFQ provides 
strong weighted fairness, and the proposed Min--Max scheduler significantly 
reduces cross-tenant slowdown imbalance while maintaining competitive throughput. 
We conclude that fairness-oriented scheduling in SSD controllers is essential 
to predictable performance in multi-tenant systems.
\end{abstract}

% ------------------------------------------------------------------
% 1. INTRODUCTION
% ------------------------------------------------------------------
\section{Introduction}

Modern SSDs operate as shared resources for multiple independent applications,
virtual machines, and containers. These tenants issue diverse read/write
patterns that contend for limited internal resources such as channels, chips,
and flash planes. Without careful scheduling, contention may lead to severe 
unfairness: one tenant may monopolize the SSD while others experience extreme 
latency slowdowns.

The default scheduling policy in MQSim, \textbf{Round Robin (RR)}, distributes
service opportunities equally across flows, but does not consider:
\begin{itemize}
    \item varying request sizes,
    \item asymmetric read and write latency,
    \item tenant-specified weights,
    \item latency-sensitive flows.
\end{itemize}
This causes unpredictable performance and poor fairness.

To address these issues, we investigate three fairness schedulers:
\begin{enumerate}
    \item \textbf{Deficit Round Robin (DRR)} — O(1), weight-based scheduling.
    \item \textbf{Quick Fair Queueing (QFQ)} — efficient weighted fair queueing.
    \item \textbf{Min--Max Fairness Scheduler (proposed)} — minimizes worst-case slowdown gap.
\end{enumerate}

All schedulers are implemented directly into MQSim’s Transaction Scheduling Unit (TSU),
and evaluated using synthetic multi-tenant workloads.

\subsection*{Contributions}
\begin{itemize}
    \item Implemented DRR, QFQ, and a novel Min--Max fairness scheduler in MQSim.
    \item Designed three multi-tenant workloads to evaluate throughput, latency, and fairness.
    \item Defined fairness metrics including Jain’s index and slowdown fairness.
    \item Showed that classic RR fails under contention, while DRR and QFQ improve fairness, and Min--Max minimizes worst-case slowdown disparity.
\end{itemize}

% ------------------------------------------------------------------
% 2. BACKGROUND
% ------------------------------------------------------------------
\section{Background}

\subsection{SSD Architecture}
SSDs are composed of multiple levels of parallelism:
\begin{itemize}
    \item host interface,
    \item request queues,
    \item channels,
    \item chips,
    \item dies and planes,
    \item flash pages and blocks.
\end{itemize}
A Transaction Scheduling Unit (TSU) selects which flow’s request is dispatched next to a channel.

\begin{figure}[h]
\centering
\fbox{\includegraphics[width=0.8\linewidth]{ssd_architecture_placeholder.pdf}}
\caption{High-level SSD architecture and TSU scheduling point (placeholder).}
\end{figure}

\subsection{Scheduling Algorithms}

\subsubsection{Round Robin (RR)}
RR cycles through flows in fixed order.  
It is simple but unfair when flows have different workloads or weights.

\subsubsection{Deficit Round Robin (DRR)}
DRR allocates each flow $i$ a quantum $Q_i$ proportional to its weight.
A deficit counter $D_i$ accumulates unused credits.

\begin{algorithm}[h]
\caption{Deficit Round Robin (DRR)}
\begin{algorithmic}[1]
\For{each flow $i$ in cyclic order}
    \State $D_i \gets D_i + Q_i$
    \While{$D_i \ge \text{size(next\_request}(i))$}
        \State schedule request from $i$
        \State $D_i \gets D_i - \text{size(request)}$
    \EndWhile
\EndFor
\end{algorithmic}
\end{algorithm}

\subsubsection{Quick Fair Queueing (QFQ)}
QFQ estimates weighted finish times without the cost of WFQ.
Requests are scheduled in increasing order of virtual finish tags.

\subsubsection{Min--Max Fairness (Proposed)}
To minimize cross-tenant slowdown disparity, we define:
\[
r_i = \frac{S_i}{w_i}
\]
where $S_i$ = total service given to flow $i$ and $w_i$ = weight.

We choose the flow with the smallest ratio $r_i$ to reduce worst-case unfairness.

\begin{algorithm}[h]
\caption{Min--Max Fairness Scheduler}
\begin{algorithmic}[1]
\State Compute $r_i = S_i / w_i$ for all flows
\State $k = \arg\min_i r_i$
\State schedule next request from flow $k$
\State update $S_k$ based on request size
\end{algorithmic}
\end{algorithm}

% ------------------------------------------------------------------
% 3. SYSTEM MODEL & WORKLOADS
% ------------------------------------------------------------------
\section{System Model and Workloads}

\subsection{Simulation Environment}

All experiments are conducted using a \textbf{custom discrete-event SSD simulator} implemented in C++17. The simulator models a multi-channel NAND flash SSD with configurable hardware parameters and pluggable scheduling policies.

\subsubsection{Device Architecture}

The SSD model implements a simplified but realistic multi-channel architecture:

\begin{itemize}
    \item \textbf{Channel Configuration}: The device supports configurable number of parallel channels (default: 8 channels, 4 chips per channel equivalent).
    \item \textbf{Bandwidth Model}: Read and write operations have asymmetric bandwidth characteristics:
    \begin{itemize}
        \item Aggregate read bandwidth: 2000 MB/s (configurable)
        \item Aggregate write bandwidth: 1200 MB/s (configurable)
        \item Per-channel bandwidth: $\text{aggregate\_BW} / \text{num\_channels}$
    \end{itemize}
    \item \textbf{Service Time Calculation}: Request service time is computed as:
    \begin{equation}
        T_{\text{service}} = \frac{\text{request\_size}}{\text{per\_channel\_bandwidth}}
    \end{equation}
    \item \textbf{Channel Scheduling}: Requests are dispatched to the first available channel using a simple first-free-channel policy. Each channel maintains independent state tracking when it becomes available.
    \item \textbf{Operation Types}: The model distinguishes between READ and WRITE operations, with writes consuming approximately 1.67$\times$ more time per byte than reads at default settings.
\end{itemize}

\subsubsection{Scheduling Policies}

The simulator implements four scheduling algorithms for comparative evaluation:

\begin{enumerate}
    \item \textbf{Round Robin (RR)}: A simple request-per-turn rotation among active users. Optimal when request sizes are uniform, as request-level fairness equals byte-level fairness.
    
    \item \textbf{Deficit Round Robin (DRR)}: Extends RR with byte-level fairness through a deficit counter mechanism. Each user receives a quantum (default: 4KB) per round. If a request fits within the accumulated deficit, it is serviced; otherwise, the deficit carries over. Supports per-user weights for proportional allocation.
    
    \item \textbf{Weighted Fair Queuing (QFQ/WFQ)}: Implements an approximation of weighted fair queuing using virtual finish time tags. Requests are tagged with virtual completion times based on their size and the user's weight. The scheduler always selects the request with the smallest virtual finish time, ensuring precise proportional fairness even under high contention.
    
    \item \textbf{Start-Gap Fair Scheduling (SGFS)}: A spatial fairness scheduler that wraps another scheduler (WFQ by default) and periodically rotates logical user IDs across SSD channels. This rotation interval (default: every 200 requests) and gap parameter (default: 1) help mitigate channel-level unfairness that can arise from static user-to-channel mappings.
\end{enumerate}

\subsubsection{Configuration Parameters}

The simulator accepts extensive configuration via command-line arguments:

\begin{itemize}
    \item \textbf{Scheduler Selection}: \texttt{--scheduler \{rr|drr|qfq|sgfs\}}
    \item \textbf{Quantum Size}: \texttt{--quantum BYTES} (default: 4096 bytes, used by DRR)
    \item \textbf{User Weights}: \texttt{--weights w1,w2,...,wn} (comma-separated, for weighted schedulers)
    \item \textbf{Channel Count}: \texttt{--channels N} (default: 8)
    \item \textbf{Bandwidth Settings}: \texttt{--read-bw MBPS} and \texttt{--write-bw MBPS}
    \item \textbf{SGFS Parameters}: \texttt{--sgfs-rotate N} and \texttt{--sgfs-gap N} for rotation interval and gap configuration
\end{itemize}

\subsubsection{Simulation Execution Model}

The simulator uses a discrete-event simulation engine with the following components:

\begin{itemize}
    \item \textbf{Event Queue}: Priority queue managing request arrivals and channel completion events, ordered by timestamp.
    \item \textbf{Request Flow}: 
    \begin{enumerate}
        \item Trace file parsing: CSV format with columns (timestamp, process\_id, type, address, size)
        \item Request admission: Arriving requests are enqueued into the scheduler
        \item User selection: Scheduler's \texttt{pick\_user()} method selects the next tenant
        \item Request dispatch: Selected request is dispatched to first available channel
        \item Completion handling: Channel completion events trigger metrics collection
    \end{enumerate}
    \item \textbf{Time Advancement}: Simulation time advances based on event timestamps, ensuring causality is preserved.
\end{itemize}

\subsection{Workloads}

Workloads are specified via CSV trace files with the following format:
\begin{verbatim}
timestamp,process_id,type,address,size
0,0,READ,1048576,4096
100,1,WRITE,2097152,8192
200,0,READ,3145728,4096
\end{verbatim}

Where:
\begin{itemize}
    \item \texttt{timestamp}: Arrival time in microseconds
    \item \texttt{process\_id}: Tenant/user identifier (0-indexed)
    \item \texttt{type}: Operation type (\texttt{READ} or \texttt{WRITE})
    \item \texttt{address}: Logical block address (used for future spatial locality modeling)
    \item \texttt{size}: Request size in bytes
\end{itemize}

\subsubsection{Scenario F1: Bandwidth Hog vs Small Tenant}

This scenario evaluates fairness when one tenant aggressively consumes bandwidth while another has modest requirements.

\begin{itemize}
    \item \textbf{Flow A (Bandwidth Hog)}: Write-heavy workload with high intensity
    \begin{itemize}
        \item Request size distribution: 8KB--64KB (mean: 32KB)
        \item Operation mix: 80\% writes, 20\% reads
        \item Inter-arrival time: 10--50$\mu$s (high arrival rate)
        \item Total requests: 10,000+
    \end{itemize}
    \item \textbf{Flow B (Small Tenant)}: Mixed read/write workload with low intensity
    \begin{itemize}
        \item Request size distribution: 4KB--8KB (fixed small sizes)
        \item Operation mix: 60\% reads, 40\% writes
        \item Inter-arrival time: 100--500$\mu$s (moderate arrival rate)
        \item Total requests: 2,000--5,000
    \end{itemize}
    \item \textbf{Evaluation Focus}: Measures whether small tenants are protected from starvation and receive fair bandwidth allocation despite the presence of an aggressive competitor.
\end{itemize}

\subsubsection{Scenario F2: Latency-Sensitive vs Background Writer}

This scenario tests scheduler behavior when latency-critical workloads compete with throughput-oriented background tasks.

\begin{itemize}
    \item \textbf{Flow 0 (Latency-Sensitive)}: Small read-dominant latency-sensitive workload
    \begin{itemize}
        \item Request size: Fixed 4KB (typical for OLTP workloads)
        \item Operation mix: 95\% reads, 5\% writes
        \item Inter-arrival time: 20--100$\mu$s (bursty, latency-critical)
        \item Latency target: P99 latency $<$ 1ms
        \item Total requests: 5,000--8,000
    \end{itemize}
    \item \textbf{Flow 1 (Background Writer)}: Large sequential writer
    \begin{itemize}
        \item Request size: 64KB--128KB (sequential patterns)
        \item Operation mix: 100\% writes
        \item Inter-arrival time: 50--200$\mu$s (sustained throughput)
        \item Total requests: 1,000--2,000
    \end{itemize}
    \item \textbf{Evaluation Focus}: Assesses whether schedulers can prioritize latency-sensitive requests while still providing reasonable throughput to background tasks. Measures tail latency (P95, P99) for Flow 0.
\end{itemize}

\subsubsection{Scenario F3: Weighted Multi-Tenant (4:2:1)}

This scenario evaluates weighted fairness across three tenants with different service level agreements (SLAs).

\begin{itemize}
    \item \textbf{Tenant Configuration}:
    \begin{itemize}
        \item Tenant 0: Weight = 4 (premium tier)
        \item Tenant 1: Weight = 2 (standard tier)
        \item Tenant 2: Weight = 1 (basic tier)
    \end{itemize}
    \item \textbf{Workload Characteristics}:
    \begin{itemize}
        \item All tenants: Mixed read/write (70\% read, 30\% write)
        \item Request sizes: 4KB--16KB (uniform distribution)
        \item Inter-arrival times: 30--150$\mu$s (similar arrival patterns)
        \item Total requests per tenant: 3,000--5,000
    \end{itemize}
    \item \textbf{Expected Allocation}: Under ideal weighted fairness, bandwidth should be allocated in a 4:2:1 ratio. Tenant 0 should receive approximately 57\% of total bandwidth, Tenant 1 receives 29\%, and Tenant 2 receives 14\%.
    \item \textbf{Evaluation Focus}: Validates that schedulers supporting weights (DRR, QFQ) achieve the target proportional allocation, while unweighted schedulers (RR) provide equal shares.
\end{itemize}

\subsubsection{Additional Workload Scenarios}

Beyond the three primary scenarios, the evaluation includes supplementary workloads to stress-test scheduler behavior:

\begin{itemize}
    \item \textbf{Uniform Load Scenario}: All tenants have identical request sizes (4KB) and arrival rates. Tests whether complex schedulers introduce unnecessary overhead when simple policies suffice.
    
    \item \textbf{Size Disparity Scenario}: Extreme request size differences (4KB vs 128KB) to evaluate byte-level fairness mechanisms in DRR and QFQ.
    
    \item \textbf{Bursty Workload}: One tenant generates bursty traffic (high arrival rate for short periods) while others maintain steady rates. Tests scheduler responsiveness to traffic pattern changes.
    
    \item \textbf{Read/Write Asymmetry}: Scenarios with varying read/write ratios across tenants to evaluate whether schedulers account for operation type differences in service time.
\end{itemize}

\subsection{Performance Metrics}

The evaluation employs comprehensive metrics to assess both fairness and performance:

\subsubsection{Throughput Metrics}

\begin{itemize}
    \item \textbf{IOPS (I/O Operations Per Second)}: Total completed requests per second, measured both per-tenant and aggregate.
    \item \textbf{Bandwidth (MB/s)}: Data transfer rate per tenant and aggregate, computed as:
    \begin{equation}
        \text{Bandwidth} = \frac{\sum \text{request\_sizes}}{T_{\text{simulation}}}
    \end{equation}
    \item \textbf{Throughput Fairness}: Ratio of achieved bandwidth to expected bandwidth based on weights or equal shares.
\end{itemize}

\subsubsection{Latency Metrics}

\begin{itemize}
    \item \textbf{Average Latency}: Mean request completion time per tenant:
    \begin{equation}
        \bar{L}_i = \frac{1}{N_i} \sum_{j=1}^{N_i} (T_{\text{finish},j} - T_{\text{arrival},j})
    \end{equation}
    where $N_i$ is the number of completed requests for tenant $i$.
    
    \item \textbf{Tail Latency}: Percentile-based latency metrics (P50, P95, P99) to capture latency distribution characteristics, critical for latency-sensitive workloads.
    
    \item \textbf{Latency Variance}: Standard deviation of latency to measure predictability and consistency.
\end{itemize}

\subsubsection{Fairness Metrics}

\begin{itemize}
    \item \textbf{Jain's Fairness Index}: Measures bandwidth fairness across tenants:
    \begin{equation}
        J(x_1, x_2, \ldots, x_n) = \frac{(\sum_{i=1}^n x_i)^2}{n \cdot \sum_{i=1}^n x_i^2}
    \end{equation}
    where $x_i$ is the total bytes transferred by tenant $i$. The index ranges from $1/n$ (worst case, one tenant dominates) to 1.0 (perfect fairness). Idle tenants are excluded from the calculation to prevent skewing.
    
    \item \textbf{Slowdown Fairness}: Compares per-tenant slowdown relative to solo-run baselines:
    \begin{equation}
        \text{Slowdown}_i = \frac{L_i^{\text{shared}}}{L_i^{\text{solo}}}
    \end{equation}
    where $L_i^{\text{shared}}$ is the average latency under shared execution and $L_i^{\text{solo}}$ is the average latency when tenant $i$ runs alone. Fairness is achieved when all tenants experience similar slowdown ratios.
    
    \item \textbf{Weighted Fairness Deviation}: For weighted scenarios, measures deviation from target allocation:
    \begin{equation}
        \text{Deviation} = \sqrt{\frac{1}{n} \sum_{i=1}^n \left(\frac{B_i}{B_{\text{total}}} - \frac{w_i}{\sum w_j}\right)^2}
    \end{equation}
    where $B_i$ is the bandwidth received by tenant $i$ and $w_i$ is its weight.
\end{itemize}

\subsubsection{Metric Collection and Reporting}

\begin{itemize}
    \item \textbf{Per-Tenant Statistics}: The simulator collects per-tenant metrics including completed request count, total bytes transferred, and average latency.
    \item \textbf{CSV Output Format}: Results are exported to CSV files with columns:
    \begin{verbatim}
    user_id,completed,avg_latency_s,total_bytes
    \end{verbatim}
    \item \textbf{Aggregate Metrics}: System-wide metrics (total throughput, overall fairness index) are computed and reported to standard output.
    \item \textbf{Post-Processing}: Python scripts (\texttt{tools/plot\_results.py}) enable visualization of fairness trends, latency distributions, and comparative analysis across schedulers.
\end{itemize}


% ------------------------------------------------------------------
% 4. IMPLEMENTATION
% ------------------------------------------------------------------
\section{Implementation in MQSim}



The following MQSim components were extended:
\begin{itemize}
    \item \texttt{TSU.cpp} --- logic for RR, DRR, QFQ, and Min--Max.
    \item \texttt{TSU.h} --- flow state structures.
    \item \texttt{ssdconfig.xml} --- added \texttt{<Scheduler>} field.
\end{itemize}

Each scheduling decision is made when a flash channel becomes available.

Solo-flow validation experiments were run to confirm correctness and baseline performance.

\section{Experimental Methodology}

\subsection{Schedulers and Scenarios}

We use MQSim as the underlying SSD simulator and evaluate four scheduling
policies implemented in its Transaction Scheduling Unit (TSU):
\begin{itemize}
    \item RR: baseline round-robin scheduling.
    \item DRR: Deficit Round Robin with per-flow quanta.
    \item QFQ: Quick Fair Queueing, an efficient WFQ approximation.
    \item MINMAX: proposed Min--Max fairness scheduler.
\end{itemize}

We design three multi-tenant scenarios, each encoded as a separate
\texttt{IO\_Scenario} block in \texttt{workload.xml}:

\begin{description}
    \item[F1: Bandwidth Hog vs Small Tenant] Two synthetic write-heavy flows
    with different average queue depths.
    \item[F2: QoS Scenario] Two synthetic read-only flows with different
    priority classes (URGENT vs HIGH).
    \item[F3: Trace-Based Scenario] A TPCC-like OLTP workload using the
    provided \texttt{tpcc-small.trace} file.
\end{description}

For each combination of scheduler and scenario, we run MQSim via:
\begin{verbatim}
./MQSim -i ssdconfig.xml -w workload.xml
\end{verbatim}
and collect the resulting \texttt{MQSim\_Results} output file.

\subsection{Metrics and Post-Processing}

From the MQSim output, we extract the following per-flow metrics:
\begin{itemize}
    \item IOPS and bandwidth,
    \item average device response time (in microseconds),
    \item total number of completed requests.
\end{itemize}

We then compute higher-level fairness metrics in a separate Python script:
\begin{itemize}
    \item Jain's fairness index over per-flow IOPS,
    \item slowdown fairness, defined as the ratio between each flow's solo-run
    throughput and its shared-run throughput,
    \item weighted fairness accuracy (for weighted scenarios).
\end{itemize}

Table~\ref{tab:exp-matrix} shows the overall experiment matrix.

\begin{table}[h]
    \centering
    \caption{Experiment matrix: schedulers vs scenarios and metrics.}
    \label{tab:exp-matrix}
    \begin{tabular}{llccc}
        \toprule
        Scenario & Description & Schedulers & Metrics & Output Files \\
        \midrule
        F1 & Hog vs small tenant &
        RR, DRR, QFQ, MINMAX &
        IOPS, latency, Jain &
        \texttt{*{\_}F1.out} \\
        F2 & QoS (read-priority) &
        RR, DRR, QFQ, MINMAX &
        IOPS, avg/99p latency &
        \texttt{*{\_}F2.out} \\
        F3 & TPCC trace &
        RR, DRR, QFQ, MINMAX &
        IOPS, Jain, slowdown &
        \texttt{*{\_}F3.out} \\
        \bottomrule
    \end{tabular}
\end{table}


% ------------------------------------------------------------------
% 5. EXPERIMENTAL RESULTS (WITH PLACEHOLDER VALUES)
% ------------------------------------------------------------------
\section{Experimental Results}

\subsection{Throughput (IOPS)}

\begin{table}[h]
\centering
\caption{Scenario F1 Throughput (IOPS). Placeholder values.}
\begin{tabular}{lccc}
\toprule
Scheduler & Flow A & Flow B & Total \\
\midrule
RR   & XXX & XXX & XXX \\
DRR  & XXX & XXX & XXX \\
QFQ  & XXX & XXX & XXX \\
MINMAX & XXX & XXX & XXX \\
\bottomrule
\end{tabular}
\end{table}

\subsection{Latency}

\begin{table}[h]
\centering
\caption{Scenario F2 Latency (\si{\micro\second}). Placeholder values.}
\begin{tabular}{lccc}
\toprule
Scheduler & Flow 0 Avg & Flow 0 99p & Flow 1 Avg \\
\midrule
RR   & XXX & XXX & XXX \\
DRR  & XXX & XXX & XXX \\
QFQ  & XXX & XXX & XXX \\
MINMAX & XXX & XXX & XXX \\
\bottomrule
\end{tabular}
\end{table}

\subsection{Jain's Fairness Index}

\begin{table}[h]
\centering
\caption{Scenario F3 Weighted Fairness.}
\begin{tabular}{lcc}
\toprule
Scheduler & Jain's Index & Weighted Accuracy \\
\midrule
RR   & XXX & XXX \\
DRR  & XXX & XXX \\
QFQ  & XXX & XXX \\
MINMAX & XXX & XXX \\
\bottomrule
\end{tabular}
\end{table}

\subsection{Slowdown Fairness}

\begin{table}[h]
\centering
\caption{Slowdown fairness across flows for Scenario F3.}
\begin{tabular}{lccc}
\toprule
Flow & RR & DRR & QFQ & Min--Max \\
\midrule
0 & XXX & XXX & XXX & XXX \\
1 & XXX & XXX & XXX & XXX \\
2 & XXX & XXX & XXX & XXX \\
\bottomrule
\end{tabular}
\end{table}

% ------------------------------------------------------------------
% 6. DISCUSSION
% ------------------------------------------------------------------
\section{Discussion}
RR exhibits severe throughput imbalance, especially when flows have different
arrival rates or write ratios. DRR significantly improves fairness while adding 
minimal computational cost. QFQ performs best under weighted workloads.  
The proposed Min--Max scheduler minimizes worst-case slowdown disparity, making it 
a strong choice for QoS-sensitive SSD environments.

% ------------------------------------------------------------------
% 7. CONCLUSION
% ------------------------------------------------------------------
\section{Conclusion}
We implemented RR, DRR, QFQ, and a novel Min--Max fairness scheduler in 
MQSim and evaluated their performance under multi-tenant workloads. 
Preliminary results indicate that fairness scheduling can greatly improve 
predictability and QoS in SSDs with minimal throughput sacrifice. 
Future work includes integrating fairness scheduling with garbage collection 
policies and extending Min--Max fairness to multi-level flash subsystems.

\end{document}
